\appendix
%

%
%
%
%
%
%
%

%

%

%
%
%
%
%
%
%
%
%
%
%

%

%

%

%

%


\section{Reading List}
The ``large'' themes
covered in the seminar, together with the primary readings, were as follows.  The themes
were not covered consecutively -- we would often interweave them from one week to the next, which, we felt, created a somewhat lighter spirit;
and although we largely present the readings in chronological order below, they were usually not covered in that way.

 

\begin{itemize}

\item {\em The emergence of  modern geometry over the 19th century (5 sessions)\/}  

We read  Lobachevsky's work~\cite{lobachevsky_parallells}  on hyperbolic geometry,   the work of Riemann~\cite{riemann_hypothesen} and
Helmholtz~\cite{helmholtz_thatsachen,helmholtz_origin,helmholtz_origin_part2},    and concluded with Klein ~\cite{klein_nicht_euklidische,klein_erlangen}. We also read Hilbert's axiomatization~\cite{hilbert_grundlagen} of Euclidean geometry together with Brouwer's  inaugural address ~\cite{brouwer_intuitionism}.
We concluding by reading Mumford~\cite{mumford_stochasticity} and  Thurston~\cite{thurston_proof_progress,thurston_weeks_three_dimensional}
for some modern viewpoints. Non-technical material included
 correspondence between Gauss and Bolyai Sr;  Poincar{\'e}'s review of Hilbert's book;
and popular writings of Einstein.

\item {\em From atomism to probability (2 session) \/} 

We traced a path  back from Brownian motions to arguments about the molecular nature of matter, reading Maxwell~\cite{maxwell_molecules,maxwell_gases} and     
 excerpts from Mach~\cite{mach_mechanics},    as well as
Einstein~\cite{einstein_brownian},  and parts of Wiener~\cite{wiener_differential_space}.
Other material included 
 Wiener’s writings on science
  and society and Perrin’s experimental verification of Einstein’s work.

 
\item {\em From laws of thought to the automation of mathematics (3 sessions) \/}  

We studied  ``classical'' theories relating thought and logic, including
Boole~\cite{boole_laws}, Leibniz~\cite{leibniz_characteristic}, Peirce~\cite{peirce_algebra_writings}, 
and material from the dawn of the computer age --  McCulloch and Pitts ~\cite{mcculloch_pitts_calculus}
and von Neumann~\cite{vonneumann_mathematician,vonneumann_automata,vonneumann_computer_brain}.
We then read Dick's sociological study~\cite{dick_aftermath}  of early automated theorem proving   
and contrasted it with recent examples chosen by Constantin Kogler:   Georgiev {\em et al}~\cite{georgiev_mathematical}, Bryan {\em et al}~\cite{bryan_motivic}, Schmitt~\cite{schmitt_extremal}, Sothanaphan~\cite{sothanaphan_erdos}.  
Other material included Turing \cite{Turing}, more Leibniz, and Hilbert's K{\"o}nigsberg address.
 

 

 
\item {\em Newton and 17th century science (4 sessions) \/}.


  We studied the first three chapters of the {\em Principia\/}~\cite{newton_principia} carefully,
  culminating in the derivation of the universal law of gravitation from Kepler's area law. 
We studied also Newton's treatment of fluxions~\cite{newton_fluxions},  which we contrasted with Leibniz \cite{leibniz_new_method},
and  Newton's early adoption, and later rejection, of Descartes' work~\cite{descartes_geometrie}.
Finally we went back to Kepler~\cite{kepler_harmonice}. 
Other material included works both about and by Newton \cite{guicciardini, keynes_newton, newton_prophecies}, modern treatments
of the Kepler problem, and
Pauli's~\cite{pauli_kepler} Jung-ian analysis of Kepler.
\end{itemize}

We also had a number of standalone sessions, not part of any planned arcs of the seminar -- although in some cases
they reflected ideas that occurred across many of the other lectures.

\begin{itemize}
\item {\em Recreational mathematics. \/}   We studied recreational mathematics through
readings curated with Evelyn Lamb: Martin Gardner's columns on tiling problems~\cite{gardner_tessellating},
  Schattschneider~\cite{Schattschneider,schattschneider_tiling} and  Wang~\cite{wang_games}.
 

\item {\em The birth of category theory.\/} 
We studied the founding papers of category theory -- work of Eilenberg, Maclane and Steenrod~\cite{eilenberg_maclane_homomorphisms,eilenberg_maclane_extensions,eilenberg_maclane_equivalences, eilenberg_steenrod_axiomatic}, alongside the collected correspondence between Steenrod and Raymond Wilder~\cite{wilder_steenrod_letters}. These readings were chosen with Alma Steingart.

\item {\em Women in mathematics.\/}
 The creation of opportunities for women in mathematics was discussed through the lens of  broader social movements.
 The readings (curated with Karen Uhlenbeck) were primarily biographical:  
 Kessel~\cite{kessel_tenured_women}, Reid~\cite{reid_julia_robinson}, Jackson~\cite{jackson_miracle_group}, Uhlenbeck~\cite{uhlenbeck_glass_ceiling}.



\item {\em Mathematics and divinity.\/}  We looked at the relation of  mathematics and theology across several  eras: from classical Greece,  
Plato~\cite{plato_epinomis}; from early Christianity, Augustine of Hippo~\cite{augustine_free_choice}, and from the 20th century,
the wartime correspondence between Andr\'{e} and Simone Weil
~\cite{weil_correspondence}. This reading was curated with David Nirenberg, whose book~\cite{nirenberg_uncountable} was also part of the reading.
%


 \end{itemize}

%
%
%
%
%
%
%
%
%
%
%
%
%
%

\section{Statement on AI usage}
The authors made use of large language models  to critique and analyze the text, as well as to assist with preparation and verification of   some items in the bibliography.
 
